%==============================================================================
\chapter{Cas d'utilisation détaillés}
%==============================================================================
Cette partie décrit comment le système réalise les opérations des utilisateurs.
Les acteurs sont le fermier, l'agent apiculteur, l'agent superviseur, le
responsable, l'administrateur et les applications tierces.

Le diagramme de cas d'utilisation global (figure \ref{fig:uc}) présente
l'ensemble des acteurs, la généralisation de l'agent superviseur à partir de
l'agent apiculteur, ainsi que toutes les dépendances entre cas (relations
d'inclusion et d'extension).
\begin{figure}[htbp]
\centering
\includegraphics[width=\linewidth]{diagramme_cas_utilisation}
\caption{Diagramme de cas d'utilisation global avec acteurs et dépendances (include et extend).}
\label{fig:uc}
\end{figure}

\section{UC1 : Planifier et approuver une visite}
\begin{itemize}
  \item \textbf{Acteur} : agent superviseur.
  \item \textbf{Préconditions} : la ruche existe et l'agent est authentifié.
  \item \textbf{Scénario nominal} : l'agent propose une date de visite, le superviseur examine le planning puis approuve, le système enregistre le statut approuvé.
  \item \textbf{Alternatives} : le superviseur refuse et demande une nouvelle date.
  \item \textbf{Postcondition} : la visite est planifiée et visible dans le calendrier.
\end{itemize}

\section{UC2 : Réaliser une visite et remplir le rapport}
\begin{itemize}
  \item \textbf{Acteur} : agent apiculteur responsable de la ruche.
  \item \textbf{Préconditions} : une visite est planifiée et approuvée.
  \item \textbf{Scénario nominal} : l'agent renseigne date, heure, durée, raison, constatations, actions prévues et effectuées, recommandations, puis évalue l'effectif, l'état de santé et la productivité de 1 à 3.
  \item \textbf{Alternatives} : saisie hors-ligne puis synchronisation différée, ajout de photos.
  \item \textbf{Postcondition} : le rapport est enregistré et la case du calendrier passe au statut exécuté.
\end{itemize}

\section{UC3 : Consulter le tableau de bord de production}
\begin{itemize}
  \item \textbf{Acteur} : fermier ou propriétaire.
  \item \textbf{Préconditions} : des mesures de production existent.
  \item \textbf{Scénario nominal} : le fermier filtre par ferme, site et période, le système affiche les ruches dont la production dépasse le seuil et propose collecte ou extension.
  \item \textbf{Postcondition} : le fermier décide d'une intervention.
\end{itemize}

\section{UC4 : Traiter une alerte sanitaire}
\begin{itemize}
  \item \textbf{Acteur} : responsable.
  \item \textbf{Préconditions} : une baisse anormale a été détectée.
  \item \textbf{Scénario nominal} : le système signale la ruche concernée, le responsable planifie une inspection imminente.
  \item \textbf{Postcondition} : une visite d'inspection est planifiée.
\end{itemize}

\section{UC5 : Interroger la quantité de miel via l'API}
\begin{itemize}
  \item \textbf{Acteur} : application tierce (Excel ou autre).
  \item \textbf{Préconditions} : l'application dispose d'un accès sécurisé au service web.
  \item \textbf{Scénario nominal} : l'application appelle \texttt{getZummHoneyActualQuantity(zummID)}, le service calcule et retourne la quantité de miel.
  \item \textbf{Postcondition} : la valeur est restituée au format du service web XML.
\end{itemize}

